% !PRINTED 712
\weberpage{712}

\weq{13}{\frac{\vartheta'_{11}\,\vartheta_{00}(v+a)}{\pi\,\vartheta_{00}(v)\,\vartheta_{11}(a)} = -2i\sum\frac{e^{(2\nu+1)\pi i a}}{q^{2\nu+1}e^{2\pi i v}+1}.}

\weq{14}{\frac{\vartheta'_{11}\,\vartheta_{01}(v+a)}{\pi\,\vartheta_{00}(v)\,\vartheta_{10}(a)} = 2\sum\frac{(-1)^{\nu}e^{(2\nu+1)\pi i a}}{q^{2\nu+1}e^{2\pi i v}+1}.}

\weq{15}{\frac{\vartheta'_{11}\,\vartheta_{10}(v+a)}{\pi\,\vartheta_{00}(v)\,\vartheta_{01}(a)} = 2\sum\frac{e^{(2\nu+1)\pi i a}\,q^{\frac{2\nu+1}{2}}\,e^{\pi i v}}{q^{2\nu+1}e^{2\pi i v}+1}.} % [source damage] the denominator's 2\nu+1 exponent and the numerator's \pi are printed from damaged type; both readings are fixed by the parallel lines

\weq{16}{\frac{\vartheta'_{11}\,\vartheta_{11}(v+a)}{\pi\,\vartheta_{00}(v)\,\vartheta_{00}(a)} = 2i\sum\frac{(-1)^{\nu}e^{(2\nu+1)\pi i a}\,q^{\frac{2\nu+1}{2}}\,e^{\pi i v}}{q^{2\nu+1}e^{2\pi i v}+1}.}
% [sic] printed 2i before the sum, but the expansion requires -2i: (15) with a+1/2 for a, and (14) with a+omega/2 for a, both give -2i; as q -> 0 the left side behaves like 4q^{1/2} sin pi(v+a), the printed right side like its negative

\begin{center}
\rule{4cm}{0.4pt}
\end{center}
